\section{Countable diagrams and their tests}
\label{sec:colors}

The construction uses colors that encode countable diagrams with
ordinal tags.  A finite tuple passes a test when its diagram and all
subsequence diagrams are valid and agree exactly on their shared
labels.  These tests will define the object condition in
Section~\ref{sec:aec}; the endpoint construction will realize them
outside a tuple-dependent countable set of coordinates.

Restrictions and ranks of finite color patterns also occur in coloring
classes \cite[Definitions~1.1--1.7]{KolesnikovLambieHanson2016}.
There the rank measures extensions of finite monochromatic diagrams.
Here ordinal tags are assigned to the elements of countable extensional
relations, and the tests require exact agreement under restriction.

Throughout the construction we work in ZFC with CH and put
\(\delta=\omega_1\).  For \(a,a'\in 2^\omega\), write \(a=^*a'\) if
\(\{n<\omega:a(n)\ne a'(n)\}\) is finite, and set
\[
                         S=2^\omega/{=^*}.
\]
The quotient here is a set of colors; its elements will later be
realized by binary sequences in structures.

\subsection{Diagrams, codes, and restrictions}

For a positive integer \(k\), write \([k]=\{1,\ldots,k\}\) and let
\[
 \mathcal L_k
   =\{(J,i):\varnothing\ne J\subseteq[k],\ i<\omega\}.
\]
The set \(\mathcal L_k\) is countably infinite.  A label records a
nonempty set of tuple positions and an integer coordinate.

\begin{definition}[Valid diagrams]\label{col:diagram}
An \emph{arity-\(k\) diagram} is a triple
\(\mathcal D=(E,R,\rho)\), where \(E,R\) are binary relations on
\(\mathcal L_k\) and \(\rho:\mathcal L_k\to\delta\).
It is \emph{valid} if the following conditions hold.
\begin{enumerate}[label=\textup{(\roman*)}]
\item
\(E\) is an equivalence relation.  The relation \(R\) and the tags
respect \(E\): whenever \(aEa'\) and \(bEb'\),
\[
 aRb\ \Longleftrightarrow\ a'Rb',
 \qquad
 \rho(a)=\rho(a').
\]
\item
On the quotient \(Q_{\mathcal D}=\mathcal L_k/E\), use \(R\) and
\(\rho\) also for the induced relation and tag function.  They satisfy
\[
                  yRx\ \Longrightarrow\ \rho(y)<\rho(x).
\]
The relation \(R\) is extensional: distinct \(x,x'\in Q_{\mathcal D}\)
have different predecessor sets
\[
 \{y\in Q_{\mathcal D}:yRx\}
 \quad\hbox{and}\quad
 \{y\in Q_{\mathcal D}:yRx'\}.
\]
\item
The labels \((\{j\},0)\), for \(j\in[k]\), lie in pairwise distinct
\(E\)-classes.
\end{enumerate}
\end{definition}

Let \(\mathfrak D_k\) be the set of all arity-\(k\) diagrams,
including the invalid ones.

\begin{lemma}[Coding]\label{col:coding}
There are injections
\[
                  \iota_k:\mathfrak D_k\longrightarrow S
                  \qquad(k\ge1)
\]
whose ranges are pairwise disjoint.
\end{lemma}

\begin{proof}
Each \(=^*\)-class is countable: its members are obtained by changing
the values of one sequence on finitely many integers.  If
\(\kappa=|S|\), choice therefore gives
\[
                 2^{\aleph_0}\le\kappa\cdot\aleph_0.
\]
The quotient cannot be countable, since then \(2^\omega\) would be
countable.  Infinite cardinal multiplication and the reverse
inequality \(\kappa\le2^{\aleph_0}\) now give
\(|S|=2^{\aleph_0}\).

For each \(k\), there are at most \(2^{\aleph_0}\) choices for each
of \(E\) and \(R\).  By CH,
\[
 \delta^{\aleph_0}
   =(2^{\aleph_0})^{\aleph_0}
   =2^{\aleph_0},
\]
so there are at most continuum many choices for \(\rho\) as well.
It follows that
\[
 \left|\coprod_{k\ge1}\mathfrak D_k\right|
       \le\aleph_0\cdot2^{\aleph_0}=|S|.
\]
An injection from this disjoint union into \(S\) gives the stated
maps.
\end{proof}

Fix such coding maps for the remainder of the proof.  A color in
the range of \(\iota_k\) uniquely decodes to an arity-\(k\) diagram.
A color outside that range is not a code of the required arity.
The coding maps are fixed set parameters in the definition of the
class, not symbols of its language.  Allowing codes for invalid
diagrams will let us assign colors at exceptional coordinates.

We next specify restriction, including its position reindexing.
If \(\varnothing\ne J\subseteq[k]\), enumerate
\(J=\{j_1<\cdots<j_m\}\) and define
\[
 \ell_J:\mathcal L_m\longrightarrow\mathcal L_k,
 \qquad
 \ell_J(U,i)=(\{j_u:u\in U\},i).
\]
For an arity-\(k\) diagram \(\mathcal D\), its restriction
\(\mathcal D\restriction J\) is the arity-\(m\) diagram obtained
by pulling back both relations and the tag function along
\(\ell_J\).  Thus it retains exactly those labels whose position
sets are contained in \(J\).  In particular, for
\(\varnothing\ne U\subseteq[m]\),
\begin{equation}\label{col:restriction-composition}
 (\mathcal D\restriction J)\restriction U
       =\mathcal D\restriction\{j_u:u\in U\}.
\end{equation}
Validity of a diagram does not by itself guarantee extensionality
of a restriction; the tests below require validity of every
constituent diagram separately.

\begin{definition}[The full subsequence test]\label{col:test}
Given \(k\ge1\) and colors
\((s_J:\varnothing\ne J\subseteq[k])\), let
\[
                     T_k\bigl((s_J)_J\bigr)
\]
mean that both of the following conditions hold.
\begin{enumerate}[label=\textup{(\roman*)}]
\item
For every nonempty \(J\subseteq[k]\), the color \(s_J\) decodes
under \(\iota_{|J|}\) to a valid diagram \(\mathcal D_J\).
\item
For every nonempty \(J\subseteq[k]\),
\[
                    \mathcal D_J
                       =\mathcal D_{[k]}\restriction J.
\]
\end{enumerate}
For an ordered tuple \(t=(x_1,\ldots,x_k)\) of distinct elements,
write \(t\restriction J=(x_{j_1},\ldots,x_{j_m})\).
Colors attached to \(t\) and its nonempty subsequences pass the
test when \(T_k((s_{t\restriction J})_J)\) holds.
\end{definition}

The order in this definition is the order of the tuple positions.
It requires no order on the set from which tuple entries are
chosen.  Equation~\eqref{col:restriction-composition} shows that
agreement with the full diagram implies agreement between every
two nested subsequences.

\begin{lemma}[Maps between quotient diagrams]\label{col:quotients}
Suppose the test in Definition~\ref{col:test} passes.  If
\(\varnothing\ne J\subseteq[k]\), then the label map \(\ell_J\)
induces an injective map
\[
                 Q_{\mathcal D_J}\longrightarrow
                 Q_{\mathcal D_{[k]}}.
\]
This map preserves tags and preserves and reflects \(R\).
The quotient maps for nested subsequences compose.
\end{lemma}

\begin{proof}
Write \(\mathcal D_J=(E_J,R_J,\rho_J)\).
The equality
\(\mathcal D_J=\mathcal D_{[k]}\restriction J\) gives, for labels
\(a,b\in\mathcal L_{|J|}\),
\[
 aE_Jb\ \Longleftrightarrow\
       \ell_J(a)E_{[k]}\ell_J(b).
\]
The forward implication makes the quotient map well-defined,
and the reverse implication makes it injective.  The analogous
equivalence for \(R\) proves preservation and reflection of the
quotient relation; equality of the restricted tag functions
proves preservation of tags.  The maps on labels compose by
their definition, so the induced quotient maps do also.
\end{proof}

The tests permit the tuple-dependent diagrams constructed in
Section~\ref{sec:endpoint}. They do not permit one fixed color \(s_j\)
for each arity \(j\) such that
\(T_k((s_{|J|})_{\varnothing\ne J\subseteq[k]})\) holds for every
\(k\ge1\). Lemma~\ref{col:no-pattern} will prove this obstruction;
it is the ingredient that later bounds models with uncountable base.
